% Part: first-order-logic
% Chapter: axiomatic-deduction
% Section: proving-things

\documentclass[../../../include/open-logic-section]{subfiles}

\begin{document}

\iftag{FOL}
      {\olfileid{fol}{axd}{pro}}
      {\olfileid{pl}{axd}{pro}}

\olsection{\usetoken{P}{derivation}కు ఉదాహరణలు}


\begin{ex}
మనం $(\lnot !D \lor !E) \lif (!D \lif !E)$ను నిరూపించాలని
అనుకుందాం. స్పష్టంగా, ఇది మన స్వీకృతాలలో దేనికీ రూపం కాదు; అందువల్ల
దీన్ని \tetoken{నిగమనం}{!!{derive}} చేయడానికి \MP{} నియమాన్ని ఉపయోగించాలి. మనకున్న
ఏకైక నియమం~MP; అది $!A$, $!A \lif !B$ ఇచ్చినప్పుడు~$!B$ను
సమర్థించడానికి అనుమతిస్తుంది. ఒక వ్యూహం ఏమిటంటే,
\olref[prp]{ax:lor3}లో $!A$ను $\lnot !D$గా, $!B$ను $!E$గా,
$!C$ను $!D \lif !E$గా తీసుకోవడం; అంటే కింది రూపాన్ని తీసుకోవడం:
\[
(\lnot !D \lif (!D \lif !E)) \lif ((!E \lif (!D \lif !E)) \lif ((\lnot
!D \lor !E) \lif (!D \lif !E))).
\]
ఎందుకు? MPని రెండుసార్లు ఉపయోగిస్తే చివరి భాగం వస్తుంది; మనకు
కావలసిందీ అదే. అంతేకాక, $\lnot !D \lif (!D \lif !E)$ అనేది
\olref[prp]{ax:lnot2} రూపమని, $!E \lif (!D \lif !E)$ అనేది
\olref[prp]{ax:lif1} రూపమని సులభంగా చూడవచ్చు. కాబట్టి మన
\tetoken{వ్యుత్పత్తి}{!!{derivation}} ఇదే:
\begin{derivation}
  1. &  $\lnot !D \lif (!D \lif !E)$ & \olref[prp]{ax:lnot2} \\
  2. & $(\lnot !D \lif (!D \lif !E)) \lif {}$\\
  &\qquad $((!E \lif (!D \lif !E)) \lif ((\lnot !D \lor !E) \lif (!D \lif !E)))$ & \olref[prp]{ax:lor3}\\
  3. & $(!E \lif (!D \lif !E)) \lif ((\lnot !D \lor !E) \lif (!D \lif !E))$ & 1, 2, \MP\\
  4. & $!E \lif (!D \lif !E)$ & \olref[prp]{ax:lif1}\\
  5. & $(\lnot !D \lor !E) \lif (!D \lif !E)$ & 3, 4, \MP
\end{derivation}
\end{ex}

\begin{ex}\ollabel{ex:identity}
$!D \lif !D$కు ఒక \tetoken{వ్యుత్పత్తిని}{!!a{derivation}} కనుగొనే ప్రయత్నం చేద్దాం.
ఇది ఒక స్వీకృత రూపం కాదు; అందువల్ల దీన్ని \tetoken{నిగమనం}{!!{derive}} చేయడానికి
\MP{}ను ఉపయోగించాలి. \olref[prp]{ax:lif1} అనేది~$!A \lif !B$
రూపంలోని స్వీకృతం; దానికి~\MPని వర్తింపజేయవచ్చు. అయితే ఇది
ఉపయోగకరంగా ఉండాలంటే, ఈ సందర్భంలో \MP{} సరైన దశగా సమర్థించే $!B$
మనకు కావలసిన~$!D \lif !D$ అయి ఉండాలి. అప్పుడు $!A$ కూడా $!D$
అయి ఉండాలి. అంటే, \olref[prp]{ax:lif1} యొక్క ఈ రూపాన్ని
పరిశీలించవచ్చు:
\[
!D \lif (!D \lif !D)
\]
\MPని వర్తింపజేయడానికి, దానికి సరిపడే రెండవ పూర్వపక్షం~$!A$ను కూడా
సమర్థించాలి. కానీ మన సందర్భంలో అది~$!D$; $!D$ను ఒక్కదాన్నే
\tetoken{నిగమనం}{!!{derive}} చేయలేం. కాబట్టి మనకు వేరొక వ్యూహం కావాలి.

కేవలం~$\lif$ ఉన్న మరొక స్వీకృతం \olref[prp]{ax:lif2}; అంటే,
\[
(!A \lif (!B \lif !C)) \lif ((!A \lif !B) \lif (!A \lif !C))
\]
\MP{}ని రెండుసార్లు వర్తింపజేస్తే చివర్లో లోపల ఉన్న నిబంధనను
పొందవచ్చు. మళ్లీ, దానర్థం \olref[prp]{ax:lif2}లో $!A \lif !C$
మనం \tetoken{నిగమనం}{!!{derive}} చేయాలని లక్ష్యంగా పెట్టుకున్న \tetoken{సూత్రం}{!!{formula}}~$!D \lif !D$
అయ్యే రూపం కావాలి. అప్పుడు
$!A$, $!C$ రెండూ~$!D$ అవుతాయి. $!A \lif (!B \lif !C)$,
$!A \lif !B$—అంటే మన సందర్భంలో $!D \lif (!B \lif !D)$,
$!D \lif !B$—రెండూ \tetoken{వ్యుత్పాదించదగిన}{!!{derivable}} అయ్యేలా~$!B$ను ఎలా ఎంచుకోవాలి?
మనం $!B$ను ఏదిగా ఎంచుకున్నా, వీటిలో మొదటిది ఇప్పటికే
\olref[prp]{ax:lif1} రూపమే. $!B$ను $(!D \lif !D)$గా తీసుకుంటే,
$!D \lif !B$ కూడా \olref[prp]{ax:lif1} యొక్క మరొక రూపం అవుతుంది.
కాబట్టి మన \tetoken{వ్యుత్పత్తి}{!!{derivation}} ఇదే:
\begin{derivation}
  1. & $!D \lif ((!D \lif !D) \lif !D)$ & \olref[prp]{ax:lif1}\\
  2. & $(!D \lif ((!D \lif !D) \lif !D)) \lif {}$\\
  & \qquad $((!D \lif (!D \lif !D)) \lif (!D \lif !D))$ & \olref[prp]{ax:lif2}\\
  3. & $(!D \lif (!D \lif !D)) \lif (!D \lif !D)$ & 1, 2, \MP\\
  4. & $!D \lif (!D \lif !D)$ & \olref[prp]{ax:lif1}\\
  5. & $!D \lif !D$ & 3, 4, \MP
\end{derivation}
\end{ex}

\begin{ex}\ollabel{ex:chain}
కొన్నిసార్లు కొన్ని ఇతర \tetoken{సూత్రాలు}{!!{formula}s}~$\Gamma$ నుంచి ఏదో ఒక
\tetoken{సూత్రానికి}{!!{formula}} ఒక \tetoken{వ్యుత్పత్తి}{!!a{derivation}} ఉందని చూపాలనుకుంటాం. ఉదాహరణకు,
$\Gamma = \{!A \lif !B, !B \lif !C\}$ నుంచి $!A \lif !C$ను
\tetoken{నిగమనం}{!!{derive}} చేయగలమని చూపుదాం.
\begin{derivation}
  1. & $!A \lif !B$ & \Hyp\\
  2. & $!B \lif !C$ & \Hyp\\
  3. & $(!B \lif !C) \lif (!A \lif (!B \lif !C))$ & \olref[prp]{ax:lif1} \\
  4. & $!A \lif (!B \lif !C)$ & 2, 3, \MP\\
  5. & $(!A \lif (!B \lif !C)) \lif {}$\\
  & \qquad $((!A \lif !B) \lif (!A \lif !C))$ & \olref[prp]{ax:lif2}\\
  6. & $((!A \lif !B) \lif (!A \lif !C))$ & 4, 5, \MP\\
  7. & $!A \lif !C$ & 1, 6, \MP
\end{derivation}
``\Hyp'' (``పరికల్పన''కు సంకేతం) అని గుర్తించిన పంక్తులు, ఆ పంక్తిలోని
\tetoken{సూత్రం}{!!{formula}} అనేది~$\Gamma$ యొక్క ఒక \tetoken{మూలకం}{!!a{element}} అని సూచిస్తాయి.
\end{ex}

\begin{prop}
  \ollabel{prop:chain} $\Gamma \Proves !A \lif !B$, $\Gamma
  \Proves !B \lif !C$ అయితే, $\Gamma \Proves !A \lif !C$.
\end{prop}

\begin{proof}
  $\Gamma \Proves !A \lif !B$, $\Gamma \Proves !B \lif !C$ అని
  అనుకుందాం. అప్పుడు~$\Gamma$ నుంచి $!A \lif !B$కు ఒక
  \tetoken{వ్యుత్పత్తి}{!!a{derivation}}, అలాగే~$\Gamma$ నుంచి $!B \lif !C$కు ఒక
  \tetoken{వ్యుత్పత్తి}{!!a{derivation}} ఉన్నాయి. వాటిని ఒకదాని తరువాత మరొకటి చేర్చి ఒకే
  \tetoken{వ్యుత్పత్తిగా}{!!{derivation}} చేయండి. ఇప్పుడు ముందరి ఉదాహరణలోని
  \tetoken{వ్యుత్పత్తిలో}{!!{derivation}} 3--7 పంక్తులను చేర్చండి. ఇది~$\Gamma$ నుంచి
  $!A \lif !C$కు ఒక \tetoken{వ్యుత్పత్తి}{!!a{derivation}}; అదే కొత్త \tetoken{వ్యుత్పత్తిలోని}{!!{derivation}}
  చివరి పంక్తి.
  పంక్తులు 4, 7ల సమర్థనల్లో, 1వ పంక్తి సూచనను $!A \lif !B$ యొక్క
  \tetoken{వ్యుత్పత్తిలోని}{!!{derivation}} చివరి పంక్తి సూచనతో, 2వ పంక్తి సూచనను
  $!B \lif !C$ యొక్క \tetoken{వ్యుత్పత్తిలోని}{!!{derivation}} చివరి పంక్తి సూచనతో
  మార్చినా అవి చెల్లుతూనే ఉంటాయని గమనించండి.
\end{proof}

\begin{prob}
స్వీకృతాల నుంచి \tetoken{వ్యుత్పత్తులను}{!!{derivation}s} చూపడం ద్వారా కిందివి చెల్లుతాయని
నిరూపించండి:
\begin{enumerate}
\item $(!A \land !B) \lif (!B \land !A)$
\item $((!A \land !B) \lif !C) \lif (!A \lif (!B \lif !C))$
\item $\lnot(!A \lor !B) \lif \lnot !A$
\end{enumerate}
\end{prob}



\end{document}
